% pdf-header.tex
% Extra styling for Pandoc-generated PDFs
% (included via --include-in-header)

\makeatletter
\@ifundefined{savenotes}{}{%%
  \let\savenotes\relax%%
  \let\endsavenotes\relax%%
}
\makeatother

\usepackage{needspace}
\usepackage{fontspec}
\usepackage{newunicodechar}
\usepackage{fancyvrb}
\usepackage{fvextra} % allow breaklines for Pandoc's Verbatim blocks
\usepackage{framed}
\fvset{breaklines=true,breakanywhere=true}

% Provide fallback Highlighting/Shaded environments if Pandoc's template omits them
\makeatletter
\@ifundefined{VerbBar}{\newcommand{\VerbBar}{|}}{}
\@ifundefined{VERB}{\newcommand{\VERB}{\Verb[commandchars=\\\{\}]}}{}
\@ifundefined{Highlighting}{%
  \DefineVerbatimEnvironment{Highlighting}{Verbatim}{commandchars=\\\{\}}
}{}
\@ifundefined{Shaded}{%
  \newenvironment{Shaded}{\begin{snugshade}}{\end{snugshade}}
}{}
 \@ifundefined{abstract}{%
   \newenvironment{abstract}{%
     \begin{quote}
     \begingroup\centering\bfseries Abstract\par\endgroup\smallskip
   }{%
     \end{quote}
   }
 }{}
\makeatother

% Reconfigure Pandoc's Highlighting environment so literal includes wrap lines
\RecustomVerbatimEnvironment{Highlighting}{Verbatim}{%
  commandchars=\\\{\},% keep Pandoc's command chars intact
  breaklines=true,% wrap long lines instead of overflowing the margin
  breakanywhere=true
}
\defaultfontfeatures{Ligatures=TeX,Scale=MatchLowercase,Renderer=Harfbuzz}

% --------------------------------------------------
% PDF mode switching: 'print' vs 'screen'
% --------------------------------------------------
\usepackage{ifthen}
\usepackage{etoolbox}

% --------------------------------------------------
% Font configuration (XeLaTeX)
% --------------------------------------------------
\setmainfont{TeX Gyre Pagella}

\newcommand{\setnoteglyphfont}{%
  \IfFontExistsTF{Noto Color Emoji}{%
    \newfontfamily\noteglyphfont{Noto Color Emoji}[Renderer=Harfbuzz,Scale=MatchLowercase]%
  }{%
    \IfFontExistsTF{Noto Emoji}{%
      \newfontfamily\noteglyphfont{Noto Emoji}[Renderer=Harfbuzz,Scale=MatchLowercase]%
    }{%
      \IfFontExistsTF{Symbola}{%
        \newfontfamily\noteglyphfont{Symbola}[Scale=MatchLowercase]%
      }{%
        \newfontfamily\noteglyphfont{DejaVu Sans}[Scale=MatchLowercase]%
      }%
    }%
  }%
}
\setnoteglyphfont

\newcommand{\emojiFallback}[1]{{\noteglyphfont\symbol{#1}}}
\newunicodechar{📺}{\emojiFallback{"1F4FA}}
\newunicodechar{🌐}{\emojiFallback{"1F310}}
\newunicodechar{📘}{\emojiFallback{"1F4D8}}
\newunicodechar{⭐}{\emojiFallback{"2B50}}

\newcommand{\sectionheadingglyph}[1]{%%
  \llap{{\noteglyphfont #1}\hspace{0.6em}}%%
}

\newcommand{\collectionheadingglyph}[1]{%%
  \llap{{\noteglyphfont #1}\hspace{0.35em}}%% tighter hang for collection-level glyphs
}

% Pandoc will set \pdfmode via --variable=pdfmode:MODE or metadata.
% We interpret that into booleans for use below. If \pdfmode is
% undefined or not 'print', we default to 'screen'.

\newboolean{isprint}
\newboolean{isscreen}

% default to screen
\setboolean{isprint}{false}
\setboolean{isscreen}{true}

% If the user specified pdfmode=print, flip the flags.
\ifdefined\pdfmode
  \ifthenelse{\equal{\pdfmode}{print}}{%%
    \setboolean{isprint}{true}%%
    \setboolean{isscreen}{false}%%
  }{}
\fi


% --------------------------------------------------
% Page layout
% --------------------------------------------------
\ifthenelse{\boolean{isprint}}{%%
  % Print-optimized: asymmetric margins for binding + mirrored layout
  \usepackage[
    inner=1.05in,
    outer=0.85in,
    top=0.95in,
    bottom=0.95in,
    bindingoffset=0.18in,
    twoside=true
  ]{geometry}%%
}{%%
  % Screen-optimized: a bit more generous margins, single-sided
  \usepackage[margin=1.1in, twoside=false]{geometry}%%
}
\usepackage{microtype} % nicer text spacing

% --------------------------------------------------
% Page style (headers/footers)
% --------------------------------------------------
\usepackage{fancyhdr}
\newcommand{\defaultpagestylename}{fancy}

% For print, use a classic running-head layout.
% For screen, keep things subdued: footer page number only,
% and no horizontal head rule.
\ifthenelse{\boolean{isprint}}{%%
  \pagestyle{fancy}
  \fancyhf{}% clear
  % Page number at bottom center
  \fancyfoot[C]{\thepage}
  % Running heads: small italic for visual separation
  % LE = Left side on Even pages, RO = Right side on Odd pages
  \fancyhead[LE,RO]{\small\itshape\nouppercase{\leftmark}}
  % Keep default \headrulewidth (visible horizontal line)
}{%%
  % Screen: very minimal, no running heads, no top rule
  \pagestyle{fancy}
  \fancyhf{}
  \fancyfoot[C]{\thepage}
  \renewcommand{\headrulewidth}{0pt}
}


% Make sure \chaptermark exists even for 'article' class
\providecommand{\chaptermark}[1]{}

% Use chapter titles for \leftmark (no section marks in header)
\renewcommand{\chaptermark}[1]{\markboth{#1}{}}
\renewcommand{\sectionmark}[1]{}

% Make sure header area is tall enough
\setlength{\headheight}{14pt}
\setlength{\headsep}{12pt}

\newcommand{\includeBannerText}{}
\newif\ifincludeBannerActive
\includeBannerActivefalse

\newcommand{\includeBannerStripe}{%%
  \begingroup
  \setlength{\fboxsep}{6pt}%%
  \colorbox{includebannerbg}{%%
    \parbox{\dimexpr\textwidth-2\fboxsep\relax}{\centering\footnotesize\sffamily\textbf{\includeBannerText}}%%
  }%%
  \endgroup
}

\fancypagestyle{include-banner}{%%
  \fancyhf{}
  \fancyhead[C]{\includeBannerStripe}
  \fancyfoot[C]{\thepage}
  \renewcommand{\headrulewidth}{0pt}
}

\newcommand{\activateIncludeBanner}[1]{%%
  \def\includeBannerText{#1}%%
  \includeBannerActivetrue%%
  \pagestyle{include-banner}%%
}

\newcommand{\deactivateIncludeBanner}{%%
  \includeBannerActivefalse%%
  \pagestyle{\defaultpagestylename}%%
  \def\includeBannerText{}%%
}

\newenvironment{includepagebanner}[1]{%%
  \clearpage%% ensure bannered inserts always start on a fresh page
  \activateIncludeBanner{#1}%%
    \thispagestyle{include-banner}%%
}{%%
  \deactivateIncludeBanner%%
}

\newenvironment{includepagebannerchapter}[1]{%%
  \activateIncludeBanner{#1}%%
}{%%
  \deactivateIncludeBanner%%
}

\makeatletter
\apptocmd{\@chapter}{\ifincludeBannerActive\thispagestyle{include-banner}\fi}{}{}
\apptocmd{\@schapter}{\ifincludeBannerActive\thispagestyle{include-banner}\fi}{}{}
\makeatother



% --------------------------------------------------
% Part title pages: vertically offset, left-aligned, no header/footer
% --------------------------------------------------
\usepackage{ragged2e} % make sure this is present near the top

\newenvironment{parttitlepage}{%
  \clearpage
  \thispagestyle{empty}% no running head / footer
  \markboth{}{}% clear left/right marks
  \vspace*{0.32\textheight}% push content down a bit
  \begingroup
  \RaggedRight% left-aligned, no justification
}{%
  \par
  \endgroup
  \clearpage
}



% --------------------------------------------------
% Section / part / chapter formatting
% --------------------------------------------------
\usepackage{titlesec}

% PARTS (your outer "Series"/major divisions)
% Shown on their own page, big and centered
\titleformat{\part}
  [display]
  {\centering\Huge\bfseries}
  {\partname~\thepart}
  {1ex}
  {}

\titlespacing*{\part}
  {0pt}   % left
  {*6}    % before
  {*4}    % after

% CHAPTERS (your main units inside parts)
% Adjusted for a clean, book-like start
\titleformat{\chapter}
  [display]
  {\normalfont\LARGE\bfseries}
  {\chaptername~\thechapter}
  {1ex}
  {}

\titlespacing*{\chapter}
  {0pt}   % left
  {*4}    % before
  {*2}    % after

% SECTIONS (for your h2 / h3 hierarchy)
\titleformat{\section}
  {\normalfont\Large\bfseries}
  {\thesection}
  {0.5em}
  {}

\titleformat{\subsection}
  {\normalfont\large\bfseries}
  {\thesubsection}
  {0.5em}
  {}

% --------------------------------------------------
% Appendices
% --------------------------------------------------
\usepackage[toc,page]{appendix}

% When Pandoc/Lua filter switches to appendix mode
% (\appendix), chapters become "Appendix A", "Appendix B", etc.
% You can tweak the label if you like:
% \renewcommand{\appendixname}{Appendix}
% \renewcommand{\appendixtocname}{Appendices}

% --------------------------------------------------
% Optional: Sidebar environment for PDF
% (we can wire your epub:type="aside" to this later)
% --------------------------------------------------
\usepackage{mdframed}
\usepackage{xcolor}


% Sidebar background differs by mode; make it clearly visible.
\ifthenelse{\boolean{isscreen}}{%%
  % Screen PDF: slightly darker gray for clear on-screen contrast
  \colorlet{sidebarbg}{gray!20}%%
  \colorlet{calloutbg}{gray!12}%%
  \colorlet{includebannerbg}{gray!25}%%
}{%%
  % Print PDF: a bit lighter, but still clearly distinct
  \colorlet{sidebarbg}{gray!12}%%
  \colorlet{calloutbg}{gray!6}%%
  \colorlet{includebannerbg}{gray!18}%%
}

% Per-environment framing (kept in dedicated inputs for easy tweaking)
\definecolor{noteboxbg}{HTML}{FFF6D5}
\definecolor{noteboxline}{HTML}{D9A300}
\definecolor{noteboxwarningbg}{HTML}{FFF1E6}
\definecolor{noteboxwarningline}{HTML}{C05621}
\definecolor{exampleboxbg}{HTML}{E8F1FF}
\definecolor{exampleboxline}{HTML}{2F80ED}
\definecolor{definitionboxbg}{HTML}{E6F7EF}
\definecolor{definitionboxline}{HTML}{1B7F4D}

% Sidebar (epub:type="aside")
\newmdenv[
  skipabove=1.2\baselineskip,
  skipbelow=1.2\baselineskip,
  linewidth=1pt,
  linecolor=gray!54,
  roundcorner=6pt,
  backgroundcolor=sidebarbg,
  innertopmargin=0.9em,
  innerbottommargin=0.9em,
  innerleftmargin=1em,
  innerrightmargin=1em,
]{sidebar}

\newenvironment{sidebarblock}{%
  \needspace{6\baselineskip}% ensure the box does not immediately break
  \begin{sidebar}%
}{%
  \end{sidebar}%
}

\AtBeginEnvironment{sidebar}{%
  \needspace{8\baselineskip}% avoid stranded headers
  \par
  \penalty500
  \colorlet{blockquote-border}{blockquote-border-sidebar}% darken nested blockquotes
}

\AtEndEnvironment{sidebar}{%
  \par
  \penalty500
  \colorlet{blockquote-border}{blockquote-border-default}% restore default blockquote accent
}

% Sidebar glyph + badges (🗒)
\newcommand{\glyphsidebar}{{\noteglyphfont 🗒}}
\newcommand{\sidebarheadingglyph}{%
  \raisebox{0.05em}{\glyphsidebar}\hspace{0.45em}%
}

\newcommand{\sidebarinlineglyph}{%
  \sidebarheadingglyph
}

\newcommand{\sidebarinlineglyphline}{%
  \par\noindent
  \raisebox{0.05em}{\glyphsidebar}\par\vspace{0.35\baselineskip}%
}

\newcommand{\sidebarbadge}{%
  \begingroup
  \par\nointerlineskip
  \noindent
  \expandafter\colorbox\expandafter{\sidebarbadgebg}{%
    \rule{0pt}{0.80\baselineskip}% tighter line height, similar to note badges
    \hspace{0.4em}\raisebox{0.1ex}{\glyphsidebar}\hspace{0.3em}%
    {\sffamily\footnotesize SIDEBAR}\hspace{0.6em}%
  }%
  \par\nointerlineskip
  \endgroup
}

% Generic box badge fallback (used by Lua when a generic .badge is requested)
% This macro will be overridden by specialized box includes if needed.
\newcommand{\boxbadge}[1]{%
  \begingroup
  \par\nointerlineskip
  \vspace{-0.25em}%
  \noindent
  \expandafter\colorbox\expandafter{\boxbadgebg}{%
    \rule{0pt}{0.9\baselineskip}%
    \hspace{0.35em}\sffamily\footnotesize #1\hspace{0.35em}%
  }%
  \par\nointerlineskip
  \endgroup
}

% Compact inline sidebar (epub:type="sidebar-inline")
\newmdenv[
  skipabove=0.6\baselineskip,
  skipbelow=0.6\baselineskip,
  linewidth=1pt,
  linecolor=gray!54,
  roundcorner=6pt,
  backgroundcolor=sidebarbg,
  innertopmargin=0.5em,
  innerbottommargin=0.5em,
  innerleftmargin=1em,
  innerrightmargin=1em,
]{sidebartight}

\AtBeginEnvironment{sidebartight}{%
  \needspace{4\baselineskip}% keep enough room above
  \colorlet{blockquote-border}{blockquote-border-sidebar}% darken nested blockquotes
}

\AtEndEnvironment{sidebartight}{%
  \colorlet{blockquote-border}{blockquote-border-default}% restore default blockquote accent
}

% Callout boxes (epub:type="callout")
\newmdenv[
  skipabove=0.8\baselineskip,
  skipbelow=0.8\baselineskip,
  linewidth=0.6pt,
  linecolor=gray!50,
  roundcorner=4pt,
  backgroundcolor=calloutbg,
  innertopmargin=0.7em,
  innerbottommargin=0.7em,
  innerleftmargin=0.9em,
  innerrightmargin=0.9em,
]{calloutbox}

% Notes (epub:type="note")

% NOTE: The block below replaces the previous implementation and then halts
% parsing with \endinput to ignore any legacy content that may remain below.

\RequirePackage{xparse}
\RequirePackage{expl3}

%--------------------------------------------------------------------
% Glyphs
%--------------------------------------------------------------------
% Glyph font assumed defined somewhere, e.g.:
%   \newfontfamily\noteglyphfont{Noto Emoji}
% or similar.

\newcommand{\glyphnote}{{\noteglyphfont ∴}}      % neutral/default note
\newcommand{\glyphidea}{{\noteglyphfont 💡}}     % note.idea
\newcommand{\glyphinsight}{{\noteglyphfont ⁂}}  % note.insight
\newcommand{\glyphwarning}{{\noteglyphfont ※}}  % note.warning (calmer than ⚠)
\newcommand{\glyphbreadcrumb}{{\noteglyphfont ↗}} % note.breadcrumb (navigational cue)
\newcommand{\glyphcontrast}{{\noteglyphfont ⇆}}  % note.contrast
\newcommand{\glyphtransition}{{\noteglyphfont ⇥}}  % note.transition
\newcommand{\glyphnavigation}{{\noteglyphfont ⇢}} % note.navigation
\newcommand{\glyphbridge}{{\noteglyphfont ⟷}}   % note.bridge
\newcommand{\glyphusage}{{\noteglyphfont ⚙}}    % note.usage
\newcommand{\glyphintegration}{{\noteglyphfont ∞}} % note.integration
\newcommand{\glyphresponsibility}{{\noteglyphfont ⚖}} % note.responsibility
\newcommand{\glyphreflexive}{{\noteglyphfont ⟲}} % note.reflexive
\newcommand{\glyphsynthesis}{{\noteglyphfont ⊕}}   % note.synthesis
\newcommand{\glyphdisclosure}{{\noteglyphfont ◊}}   % note.disclosure
\newcommand{\glyphmisuse}{{\noteglyphfont ✕}}    % note.misuse
\newcommand{\glyphdiagnostic}{{\noteglyphfont 🔍}} % note.diagnostic
\newcommand{\glyphclosing}{{\noteglyphfont ∎}}   % note.closing
\newcommand{\glyphmethodology}{{\noteglyphfont ⊢}} % note.methodology
\newcommand{\glyphcontext}{{\noteglyphfont ◉}}   % note.context
\newcommand{\glyphscope}{{\noteglyphfont ◎}}     % note.scope
\newcommand{\glyphchecklist}{{\noteglyphfont ☐}}   % note.checklist
\newcommand{\glyphprelude}{{\noteglyphfont \textcolor{noteboxpreludeline}{◆}}}   % note.prelude


%--------------------------------------------------------------------
% Presentation primitives (namespaced as \noteheader...)
%--------------------------------------------------------------------

% 1. Inline header: small glyph before the first sentence.
\newcommand{\noteheaderinline}[1]{%
  \raisebox{0.05em}{#1}\hspace{0.4em}\ignorespaces
}

% 2. Ornamental header: a standalone glyph row.
%    Using your "lightly indented asterism" style.
% Minor insight (note.insight): small ornament row
% Minor insight (note.insight): small glyph, visibly indented
\newcommand{\noteheaderornament}[1]{%
  \par\nointerlineskip
  \vspace{0.12em}
  \begingroup
    % Tighten the row height so the ornament hugs the top padding
    \setlength{\baselineskip}{0.95\baselineskip}%
    \noindent\hspace{1.05em}\smash{\small #1}%
  \par\nointerlineskip
  \endgroup
}

% 3a. Badge strip header with neutral note color.
% Major insight (note.insight.major): badge strip aligned to left text edge
\newcommand{\noteheaderbadgenormal}[1]{%
  \begingroup
  \par\nointerlineskip
  \vspace{-0.25em}% counteract the restored top padding so the badge hugs the frame
  \noindent
  \expandafter\colorbox\expandafter{\notebadgebgnormal}{%
    \rule{0pt}{.10\baselineskip}%
    \hspace{0.35em}\raisebox{0.5ex}{\large #1}\hspace{0.35em}%
  }%
  \par\nointerlineskip
  \endgroup
}

% 3c. Badge strip header with breadcrumb accent color.
\newcommand{\noteheaderbadgebreadcrumb}[1]{%
  \begingroup
  \par\nointerlineskip
  \vspace{-0.25em}
  \noindent
  \expandafter\colorbox\expandafter{\notebadgebgbreadcrumb}{%
    \rule{0pt}{.10\baselineskip}%
    \hspace{0.35em}\raisebox{0.5ex}{\large #1}\hspace{0.35em}%
  }%
  \par\nointerlineskip
  \endgroup
}

% 3d. Badge strip header with navigation accent color.
\newcommand{\noteheaderbadgenavigation}[1]{%
  \begingroup
  \par\nointerlineskip
  \vspace{-0.25em}
  \noindent
  \expandafter\colorbox\expandafter{\notebadgebgnavigation}{%
    \rule{0pt}{.10\baselineskip}%
    \hspace{0.35em}\raisebox{0.5ex}{\large #1}\hspace{0.35em}%
  }%
  \par\nointerlineskip
  \endgroup
}

% 3e. Badge strip header with bridge accent color.
\newcommand{\noteheaderbadgebridge}[1]{%
  \begingroup
  \par\nointerlineskip
  \vspace{-0.25em}
  \noindent
  \expandafter\colorbox\expandafter{\notebadgebgbridge}{%
    \rule{0pt}{.10\baselineskip}%
    \hspace{0.35em}\raisebox{0.5ex}{\large #1}\hspace{0.35em}%
  }%
  \par\nointerlineskip
  \endgroup
}

% 3f. Badge strip header with usage accent color.
\newcommand{\noteheaderbadgeusage}[1]{%
  \begingroup
  \par\nointerlineskip
  \vspace{-0.25em}
  \noindent
  \expandafter\colorbox\expandafter{\notebadgebgusage}{%
    \rule{0pt}{.10\baselineskip}%
    \hspace{0.35em}\raisebox{0.5ex}{\large #1}\hspace{0.35em}%
  }%
  \par\nointerlineskip
  \endgroup
}

% 3g. Badge strip header with integration accent color.
\newcommand{\noteheaderbadgeintegration}[1]{%
  \begingroup
  \par\nointerlineskip
  \vspace{-0.25em}
  \noindent
  \expandafter\colorbox\expandafter{\notebadgebgintegration}{%
    \rule{0pt}{.10\baselineskip}%
    \hspace{0.35em}\raisebox{0.5ex}{\large #1}\hspace{0.35em}%
  }%
  \par\nointerlineskip
  \endgroup
}

% 3h. Badge strip header with responsibility accent color.
\newcommand{\noteheaderbadgeresponsibility}[1]{%
  \begingroup
  \par\nointerlineskip
  \vspace{-0.25em}
  \noindent
  \expandafter\colorbox\expandafter{\notebadgebgresponsibility}{%
    \rule{0pt}{.10\baselineskip}%
    \hspace{0.35em}\raisebox{0.5ex}{\large #1}\hspace{0.35em}%
  }%
  \par\nointerlineskip
  \endgroup
}

% 3i. Badge strip header with reflexive accent color.
\newcommand{\noteheaderbadgereflexive}[1]{%
  \begingroup
  \par\nointerlineskip
  \vspace{-0.25em}
  \noindent
  \expandafter\colorbox\expandafter{\notebadgebgreflexive}{%
    \rule{0pt}{.10\baselineskip}%
    \hspace{0.35em}\raisebox{0.5ex}{\large #1}\hspace{0.35em}%
  }%
  \par\nointerlineskip
  \endgroup
}

% 3f. Badge strip header with diagnostic accent color.
\newcommand{\noteheaderbadgediagnostic}[1]{%
  \begingroup
  \par\nointerlineskip
  \vspace{-0.25em}
  \noindent
  \expandafter\colorbox\expandafter{\notebadgebgdiagnostic}{%
    \rule{0pt}{.10\baselineskip}%
    \hspace{0.35em}\raisebox{0.5ex}{\large #1}\hspace{0.35em}%
  }%
  \par\nointerlineskip
  \endgroup
}

% 3g. Badge strip header with synthesis accent color.
\newcommand{\noteheaderbadgesynthesis}[1]{%
  \begingroup
  \par\nointerlineskip
  \vspace{-0.25em}
  \noindent
  \expandafter\colorbox\expandafter{\notebadgebgsynthesis}{%
    \rule{0pt}{.10\baselineskip}%
    \hspace{0.35em}\raisebox{0.5ex}{\large #1}\hspace{0.35em}%
  }%
  \par\nointerlineskip
  \endgroup
}

% 3h. Badge strip header with disclosure accent color.
\newcommand{\noteheaderbadgedisclosure}[1]{%
  \begingroup
  \par\nointerlineskip
  \vspace{-0.25em}
  \noindent
  \expandafter\colorbox\expandafter{\notebadgebgdisclosure}{%
    \rule{0pt}{.10\baselineskip}%
    \hspace{0.35em}\raisebox{0.5ex}{\large #1}\hspace{0.35em}%
  }%
  \par\nointerlineskip
  \endgroup
}

% 3b. Badge strip header with warning note color.
\newcommand{\noteheaderbadgewarning}[1]{%
  \begingroup
  \par\nointerlineskip
  \vspace{-0.25em}
  \noindent
  \expandafter\colorbox\expandafter{\notebadgebgwarning}{%
    \strut\hspace{0.25em}\raisebox{0.05em}{#1}\hspace{0.25em}%
  }%
  \par\nointerlineskip
  \endgroup
}

% 3g. Badge strip header with misuse accent color (harsher warning tone).
\newcommand{\noteheaderbadgemisuse}[1]{%
  \begingroup
  \par\nointerlineskip
  \vspace{-0.25em}
  \noindent
  \expandafter\colorbox\expandafter{\notebadgebgmisuse}{%
    \strut\hspace{0.25em}\raisebox{0.05em}{#1}\hspace{0.25em}%
  }%
  \par\nointerlineskip
  \endgroup
}

% 3h. Badge strip header with closing accent (muted gray).
\newcommand{\noteheaderbadgeclosing}[1]{%
  \begingroup
  \par\nointerlineskip
  \vspace{-0.25em}
  \noindent
  \expandafter\colorbox\expandafter{\notebadgebgclosing}{%
    \rule{0pt}{.10\baselineskip}%
    \hspace{0.35em}\raisebox{0.5ex}{\large #1}\hspace{0.35em}%
  }%
  \par\nointerlineskip
  \endgroup
}

% 3i. Badge strip header with methodology accent (indigo).
\newcommand{\noteheaderbadgemethodology}[1]{%
  \begingroup
  \par\nointerlineskip
  \vspace{-0.25em}
  \noindent
  \expandafter\colorbox\expandafter{\notebadgebgmethodology}{%
    \rule{0pt}{.10\baselineskip}%
    \hspace{0.35em}\raisebox{0.5ex}{\large #1}\hspace{0.35em}%
  }%
  \par\nointerlineskip
  \endgroup
}

% 3j. Badge strip header with context accent (warm slate).
\newcommand{\noteheaderbadgecontext}[1]{%
  \begingroup
  \par\nointerlineskip
  \vspace{-0.25em}
  \noindent
  \expandafter\colorbox\expandafter{\notebadgebgcontext}{%
    \rule{0pt}{.10\baselineskip}%
    \hspace{0.35em}\raisebox{0.5ex}{\large #1}\hspace{0.35em}%
  }%
  \par\nointerlineskip
  \endgroup
}

% 3k. Badge strip header with scope accent (warm copper).
\newcommand{\noteheaderbadgescope}[1]{%
  \begingroup
  \par\nointerlineskip
  \vspace{-0.25em}
  \noindent
  \expandafter\colorbox\expandafter{\notebadgebgscope}{%
    \rule{0pt}{.10\baselineskip}%
    \hspace{0.35em}\raisebox{0.5ex}{\large #1}\hspace{0.35em}%
  }%
  \par\nointerlineskip
  \endgroup
}

% 3l. Badge strip header with checklist accent (teal).
\newcommand{\noteheaderbadgechecklist}[1]{%
  \begingroup
  \par\nointerlineskip
  \vspace{-0.25em}
  \noindent
  \expandafter\colorbox\expandafter{\notebadgebgchecklist}{%
    \rule{0pt}{.10\baselineskip}%
    \hspace{0.35em}\raisebox{0.5ex}{\large #1}\hspace{0.35em}%
  }%
  \par\nointerlineskip
  \endgroup
}


%--------------------------------------------------------------------
% Semantic header macros (map semantic types → presentation)
%--------------------------------------------------------------------

% note.idea  → inline lightbulb
\newcommand{\headernoteidea}{%
  \noteheaderinline{\glyphidea}%
}

% note.insight (minor) → ornamental asterism row
\newcommand{\headernoteinsight}{%
  \noteheaderornament{\glyphinsight}%
}

% note.insight.major (major) → badge strip with asterism
\newcommand{\headernoteinsightmajor}{%
  \noteheaderbadgenormal{\glyphinsight}%
}

% note.warning → badge strip with warning glyph and warning color
\newcommand{\headernotewarning}{%
  \noteheaderbadgewarning{\glyphwarning}%
}

% Optional: neutral note badge (for plain notes if needed)
\newcommand{\headernotenote}{%
  \noteheaderbadgenormal{\glyphnote}%
}

% note.breadcrumb → directional badge row
\newcommand{\headernotebreadcrumb}{%
  \noteheaderbadgebreadcrumb{\glyphbreadcrumb}%
}

% note.contrast → inline bidirectional arrow
\newcommand{\headernotecontrast}{%
  \noteheaderinline{\glyphcontrast}%
}

% note.transition → inline flow arrow
\newcommand{\headernotetransition}{%
  \noteheaderinline{\glyphtransition}%
}

% note.navigation → teal navigation badge
\newcommand{\headernotenavigation}{%
  \noteheaderbadgenavigation{\glyphnavigation}%
}

% note.bridge → lavender badge emphasizing linkage
\newcommand{\headernotebridge}{%
  \noteheaderbadgebridge{\glyphbridge}%
}

% note.usage → steel-blue usage badge
\newcommand{\headernoteusage}{%
  \noteheaderbadgeusage{\glyphusage}%
}

% note.integration → teal integration badge
\newcommand{\headernoteintegration}{%
  \noteheaderbadgeintegration{\glyphintegration}%
}

% note.responsibility → amber responsibility badge
\newcommand{\headernoteresponsibility}{%
  \noteheaderbadgeresponsibility{\glyphresponsibility}%
}

% note.reflexive → slate reflexive badge
\newcommand{\headernotereflexive}{%
  \noteheaderbadgereflexive{\glyphreflexive}%
}

% note.synthesis → synthesis badge
\newcommand{\headernotesynthesis}{%
  \noteheaderbadgesynthesis{\glyphsynthesis}%
}

% note.disclosure → disclosure badge
\newcommand{\headernotedisclosure}{%
  \noteheaderbadgedisclosure{\glyphdisclosure}%
}

% note.misuse → assertive red badge
\newcommand{\headernotemisuse}{%
  \noteheaderbadgemisuse{\glyphmisuse}%
}

% note.diagnostic → magnifying-glass badge
\newcommand{\headernotediagnostic}{%
  \noteheaderbadgediagnostic{\glyphdiagnostic}%
}

% note.closing → muted closing badge
\newcommand{\headernoteclosing}{%
  \noteheaderbadgeclosing{\glyphclosing}%
}

% note.methodology → indigo methodology badge
\newcommand{\headernotemethodology}{%
  \noteheaderbadgemethodology{\glyphmethodology}%
}

% note.context → warm slate context badge
\newcommand{\headernotecontext}{%
  \noteheaderbadgecontext{\glyphcontext}%
}

% note.scope → warm copper scope badge
\newcommand{\headernotescope}{%
  \noteheaderbadgescope{\glyphscope}%
}

% note.checklist → teal checklist badge
\newcommand{\headernotechecklist}{%
  \noteheaderbadgechecklist{\glyphchecklist}%
}

% note.prelude → subtle ornament (minimal decoration)
\newcommand{\headernoteprelude}{%
  \noteheaderornament{\glyphprelude}%
}


%--------------------------------------------------------------------
% Backwards-compatible macros expected by iasp-semantics.lua
%--------------------------------------------------------------------

% Inline lightbulb for note.idea
\newcommand{\notevariantideainline}{%
  \headernoteidea
}

% Minor insight header (note.insight)
\newcommand{\notevariantinsight}{%
  \headernoteinsight
}

% Major insight header (note.insight.major)
\newcommand{\notevariantinsightmajor}{%
  \headernoteinsightmajor
}

\ExplSyntaxOn
\NewDocumentCommand{\notevariant}{m}{%
  \str_case:nnF {#1} {
    {warning}   {\headernotewarning}
    {breadcrumb}{\headernotebreadcrumb}
    {contrast}  {\headernotecontrast}
    {transition}{\headernotetransition}
    {navigation}{\headernotenavigation}
    {bridge}    {\headernotebridge}
    {usage}     {\headernoteusage}
    {integration}{\headernoteintegration}
    {responsibility}{\headernoteresponsibility}
    {reflexive}{\headernotereflexive}
    {synthesis} {\headernotesynthesis}
    {disclosure}{\headernotedisclosure}
    {misuse}    {\headernotemisuse}
    {diagnostic}{\headernotediagnostic}
    {closing}   {\headernoteclosing}
    {methodology}{\headernotemethodology}
    {context}   {\headernotecontext}
    {scope}     {\headernotescope}
    {prelude}   {\headernoteprelude}
  }{\headernotenote}%
}
\ExplSyntaxOff


%--------------------------------------------------------------------
% Note box environments
%--------------------------------------------------------------------

\newmdenv[
  skipabove=0.8\baselineskip,
  skipbelow=0.8\baselineskip,
  linewidth=0.8pt,
  linecolor=noteboxplainline,
  roundcorner=4pt,
  backgroundcolor=noteboxplainbg,
  innertopmargin=0.7em,
  innerbottommargin=0.7em,
  innerleftmargin=0.95em,
  innerrightmargin=0.95em,
]{noteboxplain}

\newmdenv[
  skipabove=0.8\baselineskip,
  skipbelow=0.8\baselineskip,
  linewidth=0.8pt,
  linecolor=noteboxline,
  roundcorner=4pt,
  backgroundcolor=noteboxbg,
  innertopmargin=0.7em,
  innerbottommargin=0.7em,
  innerleftmargin=0.95em,
  innerrightmargin=0.95em,
]{notebox}

\newmdenv[
  skipabove=0.8\baselineskip,
  skipbelow=0.8\baselineskip,
  linewidth=0.8pt,
  linecolor=noteboxbreadcrumbline,
  roundcorner=4pt,
  backgroundcolor=noteboxbreadcrumbbg,
  innertopmargin=0.7em,
  innerbottommargin=0.7em,
  innerleftmargin=0.95em,
  innerrightmargin=0.95em,
]{noteboxbreadcrumb}

\newmdenv[
  skipabove=0.8\baselineskip,
  skipbelow=0.8\baselineskip,
  linewidth=0.8pt,
  linecolor=noteboxwarningline,
  roundcorner=4pt,
  backgroundcolor=noteboxwarningbg,
  innertopmargin=0.7em,
  innerbottommargin=0.7em,
  innerleftmargin=0.95em,
  innerrightmargin=0.95em,
]{noteboxwarning}

\newmdenv[
  skipabove=0.8\baselineskip,
  skipbelow=0.8\baselineskip,
  linewidth=0.8pt,
  linecolor=noteboxcontrastline,
  roundcorner=4pt,
  backgroundcolor=noteboxcontrastbg,
  innertopmargin=0.7em,
  innerbottommargin=0.7em,
  innerleftmargin=0.95em,
  innerrightmargin=0.95em,
]{noteboxcontrast}

\newmdenv[
  skipabove=0.8\baselineskip,
  skipbelow=0.8\baselineskip,
  linewidth=0.8pt,
  linecolor=noteboxtransitionline,
  roundcorner=4pt,
  backgroundcolor=noteboxtransitionbg,
  innertopmargin=0.7em,
  innerbottommargin=0.7em,
  innerleftmargin=0.95em,
  innerrightmargin=0.95em,
]{noteboxtransition}

\newmdenv[
  skipabove=0.8\baselineskip,
  skipbelow=0.8\baselineskip,
  linewidth=0.8pt,
  linecolor=noteboxnavigationline,
  roundcorner=4pt,
  backgroundcolor=noteboxnavigationbg,
  innertopmargin=0.7em,
  innerbottommargin=0.7em,
  innerleftmargin=0.95em,
  innerrightmargin=0.95em,
]{noteboxnavigation}

\newmdenv[
  skipabove=0.8\baselineskip,
  skipbelow=0.8\baselineskip,
  linewidth=0.8pt,
  linecolor=noteboxbridgeline,
  roundcorner=4pt,
  backgroundcolor=noteboxbridgebg,
  innertopmargin=0.7em,
  innerbottommargin=0.7em,
  innerleftmargin=0.95em,
  innerrightmargin=0.95em,
]{noteboxbridge}

\newmdenv[
  skipabove=0.8\baselineskip,
  skipbelow=0.8\baselineskip,
  linewidth=0.8pt,
  linecolor=noteboxusageline,
  roundcorner=4pt,
  backgroundcolor=noteboxusagebg,
  innertopmargin=0.7em,
  innerbottommargin=0.7em,
  innerleftmargin=0.95em,
  innerrightmargin=0.95em,
]{noteboxusage}

\newmdenv[
  skipabove=0.8\baselineskip,
  skipbelow=0.8\baselineskip,
  linewidth=0.8pt,
  linecolor=noteboxintegrationline,
  roundcorner=4pt,
  backgroundcolor=noteboxintegrationbg,
  innertopmargin=0.7em,
  innerbottommargin=0.7em,
  innerleftmargin=0.95em,
  innerrightmargin=0.95em,
]{noteboxintegration}

\newmdenv[
  skipabove=0.8\baselineskip,
  skipbelow=0.8\baselineskip,
  linewidth=0.8pt,
  linecolor=noteboxresponsibilityline,
  roundcorner=4pt,
  backgroundcolor=noteboxresponsibilitybg,
  innertopmargin=0.7em,
  innerbottommargin=0.7em,
  innerleftmargin=0.95em,
  innerrightmargin=0.95em,
]{noteboxresponsibility}

\newmdenv[
  skipabove=0.8\baselineskip,
  skipbelow=0.8\baselineskip,
  linewidth=0.8pt,
  linecolor=noteboxreflexiveline,
  roundcorner=4pt,
  backgroundcolor=noteboxreflexivebg,
  innertopmargin=0.7em,
  innerbottommargin=0.7em,
  innerleftmargin=0.95em,
  innerrightmargin=0.95em,
]{noteboxreflexive}

\newmdenv[
  skipabove=0.8\baselineskip,
  skipbelow=0.8\baselineskip,
  linewidth=0.8pt,
  linecolor=noteboxsynthesisline,
  roundcorner=4pt,
  backgroundcolor=noteboxsynthesisbg,
  innertopmargin=0.7em,
  innerbottommargin=0.7em,
  innerleftmargin=0.95em,
  innerrightmargin=0.95em,
]{noteboxsynthesis}

\newmdenv[
  skipabove=0.8\baselineskip,
  skipbelow=0.8\baselineskip,
  linewidth=0.8pt,
  linecolor=noteboxdisclosureline,
  roundcorner=4pt,
  backgroundcolor=noteboxdisclosurebg,
  innertopmargin=0.7em,
  innerbottommargin=0.7em,
  innerleftmargin=0.95em,
  innerrightmargin=0.95em,
]{noteboxdisclosure}

\newmdenv[
  skipabove=0.8\baselineskip,
  skipbelow=0.8\baselineskip,
  linewidth=0.8pt,
  linecolor=noteboxmisuseline,
  roundcorner=4pt,
  backgroundcolor=noteboxmisusebg,
  innertopmargin=0.7em,
  innerbottommargin=0.7em,
  innerleftmargin=0.95em,
  innerrightmargin=0.95em,
]{noteboxmisuse}

\newmdenv[
  skipabove=0.8\baselineskip,
  skipbelow=0.8\baselineskip,
  linewidth=0.8pt,
  linecolor=noteboxdiagnosticline,
  roundcorner=4pt,
  backgroundcolor=noteboxdiagnosticbg,
  innertopmargin=0.7em,
  innerbottommargin=0.7em,
  innerleftmargin=0.95em,
  innerrightmargin=0.95em,
]{noteboxdiagnostic}

\newmdenv[
  skipabove=0.8\baselineskip,
  skipbelow=0.8\baselineskip,
  linewidth=0.8pt,
  linecolor=noteboxclosingline,
  roundcorner=4pt,
  backgroundcolor=noteboxclosingbg,
  innertopmargin=0.7em,
  innerbottommargin=0.7em,
  innerleftmargin=0.95em,
  innerrightmargin=0.95em,
]{noteboxclosing}

\newmdenv[
  skipabove=0.8\baselineskip,
  skipbelow=0.8\baselineskip,
  linewidth=0.8pt,
  linecolor=noteboxmethodologyline,
  roundcorner=4pt,
  backgroundcolor=noteboxmethodologybg,
  innertopmargin=0.7em,
  innerbottommargin=0.7em,
  innerleftmargin=0.95em,
  innerrightmargin=0.95em,
]{noteboxmethodology}

\newmdenv[
  skipabove=0.8\baselineskip,
  skipbelow=0.8\baselineskip,
  linewidth=0.8pt,
  linecolor=noteboxcontextline,
  roundcorner=4pt,
  backgroundcolor=noteboxcontextbg,
  innertopmargin=0.7em,
  innerbottommargin=0.7em,
  innerleftmargin=0.95em,
  innerrightmargin=0.95em,
]{noteboxcontext}

\newmdenv[
  skipabove=0.8\baselineskip,
  skipbelow=0.8\baselineskip,
  linewidth=0.8pt,
  linecolor=noteboxscopeline,
  roundcorner=4pt,
  backgroundcolor=noteboxscopebg,
  innertopmargin=0.7em,
  innerbottommargin=0.7em,
  innerleftmargin=0.95em,
  innerrightmargin=0.95em,
]{noteboxscope}

\newmdenv[
  skipabove=0.8\baselineskip,
  skipbelow=0.8\baselineskip,
  linewidth=0.8pt,
  linecolor=noteboxchecklistline,
  roundcorner=4pt,
  backgroundcolor=noteboxchecklistbg,
  innertopmargin=0.7em,
  innerbottommargin=0.7em,
  innerleftmargin=0.95em,
  innerrightmargin=0.95em,
]{noteboxchecklist}

\newmdenv[
  skipabove=0.8\baselineskip,
  skipbelow=0.8\baselineskip,
  linewidth=0.8pt,
  linecolor=noteboxpreludeline,
  roundcorner=4pt,
  backgroundcolor=noteboxpreludebg,
  innertopmargin=0.7em,
  innerbottommargin=0.7em,
  innerleftmargin=0.95em,
  innerrightmargin=0.95em,
]{noteboxprelude}

% Examples (epub:type="example")
\newmdenv[
  skipabove=0.85\baselineskip,
  skipbelow=0.85\baselineskip,
  linewidth=0.8pt,
  linecolor=exampleboxline,
  roundcorner=4pt,
  backgroundcolor=exampleboxbg,
  innertopmargin=0.75em,
  innerbottommargin=0.75em,
  innerleftmargin=1em,
  innerrightmargin=1em,
]{examplebox}

% Definitions (epub:type="definition")
\newmdenv[
  skipabove=0.85\baselineskip,
  skipbelow=0.85\baselineskip,
  linewidth=0.8pt,
  linecolor=definitionboxline,
  roundcorner=4pt,
  backgroundcolor=definitionboxbg,
  innertopmargin=0.75em,
  innerbottommargin=0.75em,
  innerleftmargin=1em,
  innerrightmargin=1em,
]{definitionbox}

% Bridgehead blocks (epub:type="bridgehead")
% Simple framed environment used between major series

\definecolor{bridgeheadbg}{HTML}{F5F1FF}
\definecolor{bridgeheadline}{HTML}{8062B7}

\newmdenv[
  skipabove=1.4\baselineskip,
  skipbelow=1.4\baselineskip,
  linewidth=0.8pt,
  linecolor=bridgeheadline,
  roundcorner=4pt,
  backgroundcolor=bridgeheadbg,
  innertopmargin=1.2em,
  innerbottommargin=1.1em,
  innerleftmargin=1.2em,
  innerrightmargin=1.2em,
]{bridgeheadblock}


% --------------------------------------------------
% Hyperref styling (screen vs print)
% --------------------------------------------------
\usepackage[hidelinks]{hyperref}

\ifthenelse{\boolean{isscreen}}{%%
  % Screen PDF: colored links for easier navigation
  \hypersetup{
    colorlinks=true,
    linkcolor=blue,
    urlcolor=blue,
    citecolor=blue
  }
}{%%
  % Print PDF: avoid colored/boxed links
  \hypersetup{
    colorlinks=false,
    pdfborder={0 0 0}
  }
}

% --------------------------------------------------
% Epigraphs
% --------------------------------------------------
\usepackage{ragged2e}

\newenvironment{epigraph}{
  \begin{flushright}
  \begin{minipage}{0.65\textwidth}
  \itshape
}{
  \end{minipage}
  \end{flushright}
  \vspace{2em}
}

% Slightly tighter epigraph variant for inline, chapter-opening use
\newenvironment{epigraphtight}{
  \begin{flushright}
  \begin{minipage}{0.6\textwidth}
  \itshape
}{
  \end{minipage}
  \end{flushright}
  \vspace{1em}
}



% Usage (once your Lua filter or Pandoc mappings emit it):
% \begin{sidebar}
%   Sidebar text...
% \end{sidebar}
